\section{A rank-deficient cubic stratum}\label{sec:cubic}
Let $\mathfrak F$ denote the degree-three centres with
\[
 f=z(z-r)^2,\qquad g=(z-s)^3,\qquad 0<s<r<D,
\]
strict weights $\pi<\alpha<1-\pi$, strictly positive $U,V$, $\det V\ne0$,
and at least seven clocks. The first channel may have either rank. Put
$\gamma=1-\alpha$ and
\begin{equation}\label{eq:ABE}
 A=(r+3s)(2r-3s)^2,\quad B=(3D-r)s-2rD,\quad E=\gamma A-4\pi r^3.
\end{equation}
The letter $A$ in \eqref{eq:ABE} is a scalar, not the target polynomial.

\begin{theorem}[The entire cubic fibre]\label{thm:cubic-fibre}
The exact spectral fibre is a singleton precisely on
\begin{equation}\label{eq:I}
 \mathfrak I=\{\det U\ne0\}\cup\{B<0\}\cup\{E<0\}.
\end{equation}
On its complement, besides $(f,g)$ all component-polynomial pairs are,
up to permutation,
\begin{equation}\label{eq:remote-interval}
 (f,h_c),\quad h_c=(1-c)f+cg,\quad
 c_*\le c\le\min(c_D,\gamma/\pi),
\end{equation}
where
\[
 c_*=\frac{4r^3}{A},\quad b=\frac{rs}{3s-2r},\quad
 R(z)=\frac{f(z)}{g(z)},\quad c_D=\frac{R(D)}{R(D)-1}.
\]
The interval is nonempty even at equality in either constraint, and
every one of its pairs has a positive stochastic realization. Conditional
on the component pair, that realization is unique.
\end{theorem}
\begin{proof}
Lemma~\ref{lem:normalizer} fixes $q,K$ for every competitor. If $U$ is
invertible, Lemma~\ref{lem:fullrank} applies; in particular
$\det K/\det K_3=fg$ gives spectral identification directly.

If $\det U=0$, stochasticity gives $u_1=u_2=u$. The second marginal
$V\diag(\alpha,\gamma)(f,g)^{\mathsf T}$ has coefficient rank two.
Every competitor therefore has invertible second channel and its two
monic polynomials lie in the affine pencil $h_c$. Direct expansion gives
\begin{equation}\label{eq:disc}
 \disc(h_c)=c(c-1)^2(r-s)^3(Ac-4r^3),\qquad h_c(0)=-cs^3,
\end{equation}
and
\[
 \frac{R'(z)}{R(z)}=\frac{(2r-3s)z+rs}{z(z-r)(z-s)},\qquad
 A=4r^3-27s^2(r-s)<4r^3.
\]
Nonnegative roots exclude $c<0$, and the discriminant excludes $0<c<1$.
For $c>1$ there is a root in $(s,r)$ and no root in $[0,s]$. The two
other admissible roots must be above $r$, solving $R(z)=c/(c-1)>1$.
They exist only for $s>2r/3$. In that case $R$ increases to its maximum
at $b$ and then decreases to one. The double-root value is $c_*>1$;
the larger root increases from $b$ to $D$ as $c$ increases from $c_*$
to $c_D$. Thus the whole real-rooted pencil in $[0,D]$ is $\{0,1\}$,
with $[c_*,c_D]$ added exactly when $b\le D$.

If a competing pair has pencil coordinates $c_1,c_2$, its weights obey
$\alpha'_1c_1+\alpha'_2c_2=\gamma\in(0,1)$. One coordinate must therefore
be zero, and the other is one or a value $c\ge c_*$. The latter weight
is $\gamma/c$, so the floor imposes $c\le\gamma/\pi$; the other weight
is automatically at least $\alpha>\pi$. Conversely define
\[
 \beta_f=\alpha v_1+\gamma(1-1/c)v_2,\qquad
 \beta_h=\gamma v_2/c.
\]
Both vectors are positive and
$\beta_ff+\beta_hh_c=\alpha v_1f+\gamma v_2g$.
Their entry sums are the stated weights; normalizing them and keeping
first channel $(u,u)$ gives the realization. Since the polynomials are
linearly independent, coefficient functionals dual to them recover the
two rank-one coefficient matrices uniquely, hence their weights and
normalized marginals. Finally $b\le D$ is equivalent to $B\ge0$ and,
on this set, $c_*\le\gamma/\pi$ is equivalent to $E\ge0$.
This proves the entire classification, not only its open rank-one part.
\end{proof}

\subsection{Local constants and remote separation}
Let $d_{\rm par}$ be Euclidean distance in monic coefficients and
stochastic entries, minimized over the component permutation.

\begin{lemma}[A local inverse through the walls]\label{lem:local-cubic}
For every compact $K_0\subset\mathfrak F$, there exist $\rho,M,t_0>0$
such that a closed-model competitor whose component pair is within
$\rho$ of $(f,g)$, up to permutation, satisfies
\begin{equation}\label{eq:local-inverse}
 d_{\rm par}(\xi,\theta)\le M\max_j\|F_T(\xi)_j-P_{\theta,j}\|_F
\end{equation}
whenever the observation error is below $t_0$. These constants require
no lower bound on $|\det U|$, $-B$, or $-E$.
For a compact $H\subset\mathfrak I$ and a sufficiently small fixed
$\rho$, the remote separation
\begin{equation}\label{eq:separation}
 \eta_H(\rho)=\min_{\substack{\theta\in H,\ \xi\text{ in the closed model}\\
                         d_{\rm par}(\xi,\theta)\ge\rho}}
                   h_w(F_T(\xi),P_\theta)
\end{equation}
is positive whenever the minimizing set is nonempty. In particular the
whole-model inverse holds uniformly on $H$.
\end{lemma}
\begin{proof}
We make the bounds independent of the rank of $U$. Normalizer recovery
controls $q'-q,K'-K$ by $C\delta$, where $\delta$ is the maximum
observation error. Let $B_\theta$ be the coefficient matrix of the second
marginal and set $\beta=\sigma_2(B_\theta)>0$. Writing
$B'=V'\diag(\alpha')F'_{\coef}$, compactness gives a bound
$M_0\ge\|\diag(\alpha')F'_{\coef}\|$. Consequently
\[
 \beta/2\le\sigma_2(B')\le\sigma_{\min}(V')M_0,
 \qquad \|\diag(\alpha')^{-1}V'^{-1}\|
                   \le\frac{2M_0}{\pi\beta}.
\]
Multiplication by this bounded inverse expresses each competitor
polynomial as a bounded linear combination of $f,g$ with error $O(\delta)$.
Monicity makes its coefficient sum $1+O(\delta)$, giving
$f'_b=h_{c_b}+O_{\coef}(\delta)$. The prescribed small coefficient
neighbourhood makes $c_b$ close to zero and one respectively.

Near zero the endpoint inequality $f'_b(0)\le0$ gives $c_b\ge-C\delta$.
The derivative of \eqref{eq:disc} at zero is $-4r^3(r-s)^3<0$;
nonnegative discriminant gives $c_b\le C\delta$ after absorbing a
quadratic term. Near one put $v=1-c_b$ and depress the competing cubic
using its own mean. Its two depressed coefficients satisfy
\[
 a=f'(s)v+O(v^2+\delta),\qquad b_0=f(s)v+O(v^2+\delta).
\]
Real-rootedness gives $a\le0$ and
$|b_0|\le2(-a)^{3/2}/(3\sqrt3)$. Since $f(s)=s(r-s)^2>0$,
\[
 |v|\le C\delta+Cv^2+C(|v|+\delta)^{3/2}.
\]
Choose the coefficient radius small enough to absorb the last terms.
This yields $|v|\le C\delta$. Fixed coefficient functionals dual to
$f,g$ now have an $I+O(\delta)$ value matrix on the competing pair.
Applied to $K'$, they recover the two coefficient matrices with
$O(\delta)$ error. Entry sums recover weights and normalized row and
column sums recover both channels. This proves \eqref{eq:local-inverse}.

Here is a precise list of the uniform margins used in this argument:
the smallest singular value of the normalization operator $N_\theta$,
$\sigma_2(B_\theta)$, the Gram determinant of the coefficient pair
$(f,g)$, $4r^3(r-s)^3$, $s(r-s)^2$, the positive weight floor, and
the root and positive-channel margins on $K_0$. Their minima are positive
and all upper coefficient and Taylor bounds are finite. Choose $\rho$
so that the displayed absorption constants work with these common
bounds, then $t_0$ so that $\sigma_2(B')\ge\beta/2$. This supplies a
single $M,\rho,t_0$, rather than a rank-two inverse whose constant
could diverge at $\det U=0$.

For \eqref{eq:separation}, the feasible pairs $(\theta,\xi)$ form a
compact set. A zero minimum would yield an exact competitor separated
from the centre in parameter distance. Theorem~\ref{thm:cubic-fibre}
and its uniqueness statement exclude this on $H$. A sufficiently small
observation ball therefore lies in the local coefficient neighbourhood,
where the first part applies. This last exclusion uses identification;
the local inverse does not. In particular $\eta_H(\rho)$, not the local
coefficient constant, is the quantity allowed to degenerate at a remote
identification wall.
\end{proof}

\subsection{The complete local Fisher fibre}
Use variables $a,X,Y\ge0$, free $b,c$, and five free stochastic
coordinates $\eta$. The component jets are
\begin{equation}\label{eq:cubic-jets}
 p_f=-a(z-r)^2-bz(z-r)-Xz,\qquad
 p_g=-c(z-s)^2-Y(z-s).
\end{equation}
Let $L$ be the seven free score columns and $D_0$ the three columns
indexed by $0,x,y$ for $(a,X,Y)$. Define $Q$ by \eqref{eq:mult-Q} and
\begin{equation}\label{eq:cubic-kappa}
 J(X,Y)=\min_{a\ge0}(a,X,Y)Q(a,X,Y)^{\mathsf T},\quad
 \kappa_x=\min_{Y\ge0}J(1,Y),\quad
 \kappa_y=\min_{X\ge0}J(X,1).
\end{equation}

\begin{proposition}[Coercivity and finite supports]\label{prop:cubic-Q}
Throughout $\mathfrak F$, $L$ has rank seven, $Q_{00}>0$, and
$\kappa_x,\kappa_y$ are positive and locally Lipschitz. For every compact
$K_0\subset\mathfrak F$,
\[
 \|S_\theta(a,b,c,\eta,X,Y)\|_{I,\theta}
       \ge c_{K_0}\|(a,b,c,\eta,X,Y)\|\quad(a,X,Y\ge0).
\]
Each $\kappa_i$ is given by four finite support tests. For
$A\subseteq\{0,x,y\}\setminus\{i\}$ with $Q_{AA}$ positive definite, put
\[
 v_A=-Q_{AA}^{-1}Q_{Ai},\qquad k_{i,A}=Q_{ii}-Q_{iA}Q_{AA}^{-1}Q_{Ai}.
\]
The test is $v_A\ge0$ and $Q_{ji}+Q_{jA}v_A\ge0$ for inactive $j$.
The empty support has value $Q_{ii}$. At least one test succeeds, and
all successful tests give $\kappa_i$.
\end{proposition}
\begin{proof}
Every feasible jet is realized by roots
\[
 ta,\qquad r+tb/2\pm\sqrt{tX},\qquad
 s+tc/3+\sqrt t\,y_i,\quad \sum_i y_i=0,\quad\tfrac12\sum_i y_i^2=Y,
\]
and order-$t$ stochastic increments. Their observation increment is
$tS+O(t^{3/2})$. A zero score, combined with the local inverse of
Lemma~\ref{lem:local-cubic}, forces both component jets and $\eta$ to
vanish. Evaluation and coefficient comparison in \eqref{eq:cubic-jets}
then force every variable to be zero. Compactness of the feasible unit
sphere gives coercivity, uniformly on $K_0$. It also proves independence
of the free columns, $Q_{00}>0$, and existence, positivity and boundedness
of the constrained minimizers.

Choose a minimizer with minimal positive support in the projected
score columns. Those supported columns are independent: a nonzero
linear dependence would allow a signed change of the positive
coefficients, preserving the represented vector, until one coefficient
vanished. On an independent support the stationarity and inactive
inequalities are exactly the stated tests. Conversely their KKT
conditions prove optimality by convexity. This handles a singular full
$Q$ without inverting it. Uniformly bounded minimizing variables and
smooth dependence of the score show local Lipschitz dependence of the
minimum by evaluating each problem at a minimizer of the other.
\end{proof}

\begin{theorem}[Local normal form and whole-model restriction]\label{thm:cubic-local}
For every centre of $\mathfrak F$ the local leading root set is
\begin{equation}\label{eq:cubic-leading}
 \bigl(\{0\},\{-\sqrt X,\sqrt X\},Z(z^3-Yz-Z)\bigr),\qquad
 X,Y\ge0,\quad J(X,Y)\le4,\quad27Z^2\le4Y^3.
\end{equation}
On compact subsets of $\mathfrak F$, scaled local observation balls
converge to this set with Hausdorff error $O(\sqrt t)$. Hence
\begin{equation}\label{eq:cubic-C}
 \omega_{\theta,\loc}(t)=C_{\loc}(\theta)\sqrt t+O(t),\qquad
 C_{\loc}=\max\{\sqrt2\,\kappa_x^{-1/4},\ 2\sqrt{2/3}\,\kappa_y^{-1/4}\}.
\end{equation}
The error is compact-uniform, and $C_{\loc}$ is positive and locally
Lipschitz on all of $\mathfrak F$. On compact subsets of $\mathfrak I$
these statements hold for the whole-model modulus. The support tests,
\eqref{eq:I}, and $9\kappa_y-16\kappa_x$ give a finite semialgebraic
partition. On its constant choices,
\[
 C_{\loc}^4=\max\{4/\kappa_x,64/(9\kappa_y)\}
\]
is one rational function, and $C_{\loc}$ is Nash after refinement.
\end{theorem}
\begin{proof}
The local inverse bounds coefficients and stochastic increments by
$O(t)$ for every local competitor. Its first component has exact roots
$ta,r+tb/2\pm\sqrt{tX}$ with bounded $a,b,X$ and $a,X\ge0$.
Its second component has exact roots $s+tc/3+\sqrt t\,y_i$ with
$\sum_i y_i=0$ and bounded squared sum. Put
$Y=\tfrac12\sum_i y_i^2$ and $Z=\prod_i y_i$. The real-cubic condition
is exactly $Y\ge0,27Z^2\le4Y^3$. Expansion gives
$F=P+tS+O(t^{3/2})$ and $h_w=t\|S\|_I/2+O(t^{3/2})$ uniformly on
these bounded coordinates. The radial contraction of all first-order
variables, together with $Z\mapsto\rho^{3/2}Z$, proves the outer
Hausdorff inclusion. Conversely choose bounded minimizing free and
endpoint variables and realize the root arcs in the proof of
Proposition~\ref{prop:cubic-Q}, using $\rho=1-M\sqrt t$. They lie in
the observation ball and give the inner inclusion. Thus the inequality
boundary and all joint cubic coefficient relations are realized, not
only contained in an outer tangent cone.

Homogeneity gives $X_{\max}=2/\sqrt{\kappa_x}$ and
$Y_{\max}=2/\sqrt{\kappa_y}$. The double-cluster diameter is
$\sqrt{X_{\max}}$. Lemma~\ref{lem:ordered} with $m=3$ gives the triple
diameter $2\sqrt{Y_{\max}/3}$, also attained by
$(-2h,h,h)$ and $(-h,-h,2h)$ with $h=\sqrt{Y_{\max}/3}$.
The maximum of marginal diameters yields \eqref{eq:cubic-C}.
The remaining assertions follow from Proposition~\ref{prop:cubic-Q},
cluster separation, and the positive remote exclusion in
Lemma~\ref{lem:local-cubic}. In particular the local formula remains
valid at $B=0$ or $E=0$, but its identification with the whole-model
modulus does not.
\end{proof}

\begin{example}[The earlier rational phase values]
For $D=1$, $\pi=1/10$, $r=3/5$, $s=1/5$, $\alpha=2/5$,
$u_1=u_2=(1/3,2/3)^{\mathsf T}$,
$v_1=(1/4,3/4)^{\mathsf T}$, $v_2=(3/4,1/4)^{\mathsf T}$,
clocks $(6/5,7/5,8/5,9/5,2,5/2,3)$ and equal weights,
$B=-18/25$. The empty supports are admissible, the double cluster
dominates, and exact rational evaluation gives
\[
 C^4=\frac{78550423766145732273576408862601811503774391446}
            {385279214348552940073664077156875},\qquad
 3778.707<C<3778.708.
\]
Changing only $u_2$ to $(2/5,3/5)^{\mathsf T}$ gives the triple-dominant
phase, $59.970<C<59.971$. The exact arithmetic records for both values
are retained in the archival companion. The next section evaluates a
wall datum and its different entrance constant, rather than interpreting
these large values as evidence of an unspecified wall blow-up.
\end{example}

\begin{corollary}[Centre-known local risk]\label{cor:risk}
Let $H\subset\mathfrak I$ be compact, fix $0<c\le1/(8\sqrt2)$, and
observe independent categorical samples at the design clocks with
$n_j/N\to w_j$ uniformly on $H$ (integer rounding suffices).
For squared spectral loss on the specified neighbourhood
$h_w(F_T(\xi),P_\theta)\le cN^{-1/2}$, with the centre known, the
minimax risk $\mathcal R_N(\theta,c)$ satisfies uniformly on $H$
\[
 \frac{3c}{32}C_{\loc}(\theta)^2-O_H(N^{-1/4})
 \le\sqrt N\,\mathcal R_N(\theta,c)
 \le cC_{\loc}(\theta)^2+O_H(N^{-1/4}).
\]
\end{corollary}
\begin{proof}
A constant decision in the specified spectral ball has loss at most its
squared diameter. For a diameter pair in the ball, product Hellinger
distance is at most $1/4+O_H(N^{-1})$ with integer rounding. Total variation
is at most this distance. A two-point test gives squared-risk lower bound
$(3/32-O_H(N^{-1}))$ times the squared diameter. Substitute \eqref{eq:cubic-C}. This is an oracle local
statement; it makes no assertion of unknown-centre adaptation or honest
confidence sets across an identification wall.
\end{proof}
